\section{2교시 (90분) — EAC · Earned Schedule · 착시 · 흐름 지표 · 임계경로 회복}

% =====================================================================
% 2교시 — EAC · Earned Schedule · 착시 · 흐름 지표 · 임계경로 회복
% =====================================================================
\begin{frame}{2교시에서 배우는 것 — "이대로 가면 어디에 도착하는가"}
\begin{itemize}
\item \textbf{EAC 4공식}은 4개의 전제 — 공식 EAC는 근거가 있는 것 하나(EAC4 150.86), 검산은 EAC1$'$, 경고는 EAC3$'$
\item \textbf{TCPI}는 분모를 적어야 숫자다 — BAC / 조정 AC / 집행 한도 / EAC4
\item \textbf{Earned Schedule} — SPI는 종료로 수렴한다. 시간으로 재면 SV(t) $-$0.72개월, IEAC(t) 18.48개월
\item \textbf{착시 5개}에 이름·크기·조정값을 붙여 병기한다 — 대체하지 않는다
\item \textbf{흐름 지표}는 EVM의 2주 앞 신호, \textbf{CPM}은 청진기 — A10 $-$10이 EX-01을 만든다
\end{itemize}
\keymsg{온도계(EVM) \ra{} 보정표(ES) \ra{} 청진기(CPM) — 세 도구가 같은 환자를 다르게 말할 때 어느 것을 믿는가}
\note{교재 제2장 2.3 후반\til{}2.9, 워크북 §2(⑫ 성과보고) 항목 4\til{}9. 1교시가 "지금 어디 있는가"(편차·지수)였다면 2교시는 "이대로 가면 어디에 도착하는가"(예측)와 "왜 지표가 실제와 다르게 보이는가"(착시)이고, 마지막에 "그래서 누가 무엇을 결정하는가"(EX-01)로 끝난다. 모든 숫자는 워크북 완성 예시본 — 4교시 손계산에서 학생이 같은 값을 다시 만들게 된다. 도구 비유(온도계·보정표·청진기)를 첫 프레임에서 세워 두고 교시 내내 반복한다. 예상 소요 3분.}
\end{frame}

\begin{frame}{EAC 4공식 — 어느 공식이 맞는가가 아니라 어느 전제가 맞는가}
\begin{tbl}[\scriptsize]{L{0.30\textwidth}L{0.30\textwidth}rrL{0.18\textwidth}}
\textbf{공식} & \textbf{전제} & \textbf{EAC} & \textbf{VAC} & \textbf{판정} \\ \midrule
EAC1 = BAC ÷ CPI & 누적 원가 효율 유지 & 141.23 & +4.57 & 기각 — 착시 포함 \\
EAC2 = AC + (BAC $-$ EV) & 편차는 일회성 & 143.50 & +2.30 & 기각 — 4.70 미반영 \\
EAC3 = AC + (BAC $-$ EV) ÷ (CPI×SPI) & 원가·일정 효율 동시 & 145.43 & +0.37 & 참고 \\
\textbf{EAC4 = AC + 상향식 ETC} & 항목별 잔여 재추정 & \textbf{150.86} & \textbf{$-$5.06} & \textbf{공식} \\
EAC1$'$ = BAC ÷ 조정 CPI 0.968 & 발생주의 조정 후 유지 & 150.57 & $-$4.77 & EAC4 검산 \\
EAC3$'$ = 조정 AC + 잔여 ÷ (0.968×0.943) & 조정 CPI × SPI & \textbf{155.04} & $-$9.24 & \textbf{경고 대역} \\
\end{tbl}
\keymsg{PMB 145.8 아래의 EAC1은 고정가 SI 68.2가 줄지 않는 한 구조적으로 불가능하다 — 공식이 아니라 전제를 기각한다}
\note{§2.3 "EAC 4공식과 각각의 가정", 워크북 ⑫ 항목 5 표. 네 공식을 다 적는 이유는 보고서 안에서 전제를 논증하기 위해서다 — EAC 하나만 적는 성과보고는 대개 가장 낙관적인 EAC1을 고르고 스폰서는 원가가 남는다고 믿는다. EAC1 141.23은 CPI 1.032(라이선스 이연 착시 포함)의 외삽이고, EAC2 143.50은 미발생 4.70과 컨틴전시 잔여 집행이 잔여에 없다. EAC3$'$ 155.04는 일정 회복 실패 시의 도착지로 집행 한도까지 0.96 — 스폰서가 기억해야 할 두 번째 숫자(첫 번째는 A10 $-$10). 4교시에서 학생은 이 표를 손으로 다시 만든다. 예상 소요 6분.}
\end{frame}

\begin{frame}{EAC4의 상향식 ETC 79.72 — 행마다 근거가 있어야 "공식"이 된다}
\begin{itemize}
\item SI 잔여 계약 38.51 + 변경 대가 1.02 + 대기 잔여 0.18 — 근거: 계약 잔액·변경 로그
\item 상용SW 5.90 = 잔여 PV 1.20 + \textbf{미발생 4.70} — 근거: 지급 조건 50/50, 2027-01 통합시험 통과 시 발생
\item 인프라 7.92 = 7.40 × 1.07 — 근거: 클라우드 실사용 추세(IS-10)
\item 단말 8.90 · 전환 5.10 · 시험·감리 7.40 · PMO 4.80 — 근거: 견적·잔여 PV
\item 검산: 조정 CPI의 EAC1$'$ 150.57과 EAC4 150.86이 \textbf{0.3 안}에서 만난다 — 두 방법이 크게 다르면 어느 전제가 틀린 것
\end{itemize}
\keymsg{수행사 견적은 계약 잔액이지 발주사의 잔여 원가가 아니다 — 상향식 ETC를 견적으로 대체하지 않는다}
\note{§2.3 같은 절. 상향식 ETC의 각 행에 근거 열이 있다는 것이 EAC4가 "공식"이 되는 조건이다. 상용SW 행의 4.70이 여기 들어가므로 EAC4는 착시 1을 이미 걷어 낸 값이고, 그래서 착시를 걷어 낸 EAC1$'$과 만난다 — 만난다는 것 자체가 검산이다. 인프라 1.07은 월 0.4 \ra{} 0.5 실사용 추세를 잔여 7개월에 적용한 것. 실무 오류 — 수행사 견적으로 ETC를 대체하는 것(발주사 관점 잔여 원가에는 단말·인프라·PMO가 있다), 변경 대가 1.02를 ETC에 넣지 않는 것(CR-08\til{}11 미승인이라도 전망에는 넣고 "미반영 변경"으로 표시). 예상 소요 4분.}
\end{frame}

\begin{frame}{TCPI 4종과 VAC — 분모를 적지 않은 TCPI는 숫자가 아니다}
\begin{tbl}[\scriptsize]{L{0.30\textwidth}L{0.28\textwidth}rL{0.32\textwidth}}
\textbf{TCPI} & \textbf{산식} & \textbf{값} & \textbf{읽기} \\ \midrule
TCPI(BAC) & 72.37 ÷ 74.67 & 0.969 & 본값 — 착시 포함, 의미 없음 \\
TCPI(BAC, 조정 AC 75.83) & 72.37 ÷ 69.97 & \textbf{1.034} & 잔여 3.4\% 절감 필요 \ra{} PMB 초과 확정 \\
TCPI(집행 한도 156.00) & 72.37 ÷ 84.87 & 0.853 & 집행 한도 안은 가능 — 컨틴전시가 완충 \\
TCPI(EAC4 150.86) & 72.37 ÷ 79.73 & 0.908 & 공식 EAC 달성에 필요한 효율 \\
\end{tbl}
\begin{itemize}
\item VAC = BAC $-$ EAC4 = $-$5.06 · 집행 한도 여유 5.14 \ra{} \textbf{원가 예외 보고 불요}
\item 단 EAC3$'$ 155.04(여유 0.96) — 10-31 데이터에서 EAC4 > 153.0이면 예외 보고 준비를 \textbf{예고}(경고 A)
\item 승인 예산 168.0과 비교한 여유 17.14는 관리예비비 12.0을 포함한 남의 돈
\end{itemize}
\keymsg{"TCPI > 1.1이면 회복 불가"를 인용할 때 분모가 BAC인지 집행 한도인지 먼저 — 온담은 BAC로는 초과 확정, 한도로는 여유}
\note{§2.3 "TCPI와 VAC", 워크북 ⑫ 항목 5 하단 표. 1교시의 BAC 논쟁(145.8 vs 156.0)이 여기서 실제 결과로 갈린다 — 같은 잔여 작업 72.37을 어느 돈으로 끝내는가에 따라 1.034와 0.853. 원가 RAG가 R이 아니라 A인 이유가 이 두 줄이다(집행 한도 기준 여유 5.14, 그러나 EAC3$'$ 경고). 예고 문장("다음 보고에서 153.0을 넘으면 예외 보고를 준비한다")은 PRINCE2의 경고 등급 운용 — 예외 보고를 갑자기 내지 않기 위한 장치다. 실무 오류: VAC를 승인 예산과 비교하는 것, CV 양수를 절감으로 읽는 것 — 설명할 수 없는 양수 CV는 대개 AC 인식의 지연. 예상 소요 5분.}
\end{frame}

\begin{frame}{EAC 사다리 — 스폰서가 기억할 숫자는 둘이다}
\fig[0.58]{fig_2_3_eac_ladder}
\figcap{그림 2-3. 온담 2026-09-30의 EAC 사다리. 네 공식은 네 전제이고, 항목별 근거를 가진 EAC4(150.86)만 공식이 될 수 있으며 조정 CPI의 EAC1$'$(150.57)이 검산한다. EAC3$'$(155.04)은 회복 실패 시 도착지 — 집행 한도까지 0.96}
\keymsg{승인 168.0 / 집행 한도 156.0 / PMB 145.8 세 층 사이에 EAC 여섯 값을 놓으면 "원가는 A"가 한눈에 읽힌다}
\note{§2.3 그림 2-3. 사다리를 아래에서 위로 읽는다 — EAC1 141.23(PMB 아래, 불가능) \ra{} EAC2 \ra{} EAC3 \ra{} PMB 145.8 \ra{} EAC1$'$·EAC4(150.6\til{}150.9, 공식) \ra{} EAC3$'$ 155.04 \ra{} 집행 한도 156.0 \ra{} 승인 168.0. 오른쪽 TCPI 4값이 각 분모에 대응한다. 학생에게 묻기 — "스폰서 정미란에게 한 문장으로 원가를 보고하라": 정답은 "공식 EAC 150.86으로 PMB를 5.06 넘지만 집행 한도 안(여유 5.14)이며, 일정 회복이 실패하면 155.04까지 간다"이다. 예상 소요 3분.}
\end{frame}

\begin{frame}{SPI의 결함 — 6개월 늦게 끝난 프로젝트의 마지막 SPI도 1.0이다}
\begin{itemize}
\item EV는 결국 BAC에 도달하고 PV도 계획 종료 후 BAC에 머문다 \ra{} SPI = EV ÷ PV는 \textbf{종료 시점에 반드시 1.0}
\item 수렴은 종료 순간이 아니라 후반부 내내 진행 — 마지막 1/3에서 SPI는 실제 지연과 무관하게 \textbf{좋아지는 것처럼} 보인다
\item SV(돈 단위)도 같은 이유로 0으로 수렴 — Lipke(2003) 「Schedule is Different」
\item 해법: 일정 편차를 돈이 아니라 \textbf{시간}으로 잰다 — Earned Schedule
\item ES = "현재 누적 EV를 PV 곡선이 계획상 \textbf{언제} 달성하기로 했는가"의 시점
\end{itemize}
\keymsg{후반부 프로젝트에서 SPI를 보고 안심하는 조직이 겪는 일 — 온담은 진척 50\%에서 이미 벌어지기 시작했다}
\note{§2.4 "결함". 칠판에 EV·PV 두 곡선을 그려 종료 직전 두 곡선이 같은 BAC 수평선에 붙는 모습을 보이면 결함이 한 번에 이해된다. 학생 대부분이 SPI를 "일정 효율"로만 알고 있고, 종료 수렴을 처음 듣는다 — 여기서 "여러분 조직의 마지막 월간 보고서 SPI가 얼마였는가"를 묻는다. 대개 1.0 근처이고 그 프로젝트가 정시에 끝났는지 물으면 대답이 갈린다. 예상 소요 4분.}
\end{frame}

\begin{frame}{Earned Schedule — ES · SV(t) · SPI(t) · IEAC(t), 온담 2026-09-30}
\begin{tbl}[\scriptsize]{L{0.22\textwidth}L{0.40\textwidth}L{0.30\textwidth}}
\textbf{지표} & \textbf{산식 · 온담 값} & \textbf{읽기} \\ \midrule
ES & C + (EV $-$ PV$_C$) ÷ (PV$_{C+1}$ $-$ PV$_C$) = 8 + 1.69 ÷ 6.09 = \textbf{8.28개월} & 8월 71.74 < EV 73.43 < 9월 77.83 \ra{} C = 8 \\
SV(t) & ES $-$ AT = 8.28 $-$ 9 = \textbf{$-$0.72개월} & ≈ $-$15영업일(평균) \\
SPI(t) & ES ÷ AT = \textbf{0.920} & SPI 0.943과 0.023 차이 \\
IEAC(t) & PD ÷ SPI(t) = 17 ÷ 0.920 = \textbf{18.48개월} & +1.48개월 ≈ +31영업일 \\
추세 종료 & \textbf{2027-07-12} vs 허용범위 06-18(+15영업일) & \textbf{초과(추세)} \\
\end{tbl}
\begin{itemize}
\item IEAC(t)는 추세 외삽이지 계획이 아니다 — 회복 계획(to-be MG3 02-19)과 나란히 놓고 어느 것이 계획인지 적는다
\item ES는 프로젝트 전체의 \textbf{평균} — 지연은 임계경로 A10에 $-$10으로 몰려 있다
\end{itemize}
\keymsg{ES는 온도계의 눈금을 바로잡는 보정표이지 청진기가 아니다 — 청진기는 CPM}
\note{§2.4 "Earned Schedule", 워크북 ⑫ 항목 6. 보간 계산을 학생과 함께 한 번 한다 — PV 누적 표(Day2 §4.5)에서 EV 73.43을 넘지 않는 마지막 달(8월 71.74)을 찾고 9월 증분 6.09로 나눈다. AT = 9(1월 착수, 9월 말). PD 17은 기준선 계획 기간(2027-05 종료). +31영업일은 1.48개월 × 약 21영업일. 실무 오류 — SPI만 보고 "회복 중"이라고 읽는 것, ES 없이 SV를 일수로 환산하는 것(1교시 봉우리 문제), IEAC(t)를 종료일로 보고하는 것, ES를 계산하고 CPM을 보지 않는 것. 워크북 항목 6의 추세 도착일 표기가 두 곳(07-12 / 07-14)으로 다르다 — 성과보고 요약표의 07-12를 정본으로 쓰고 4교시에서 정정한다. 예상 소요 7분.}
\end{frame}

\begin{frame}{SPI vs SPI(t) 월별 — 7월부터 벌어진다}
\begin{tbl}[\scriptsize]{lrrrL{0.46\textwidth}}
\textbf{월} & \textbf{SPI} & \textbf{SPI(t)} & \textbf{차이} & \textbf{판독} \\ \midrule
2026-04 & 0.784 & 0.887 & +0.103 & 조달 이연 — 돈으로 $-$8.0, 시간으로 $-$0.45개월 \\
2026-05 & 0.871 & 0.868 & $-$0.003 & 라이선스 인도로 SPI 회복, SPI(t)는 R-05 반영 \\
2026-07 & 0.914 & 0.904 & $-$0.010 & \textbf{벌어지기 시작} \\
2026-08 & 0.913 & 0.914 & +0.001 & \\
\textbf{2026-09} & \textbf{0.943} & \textbf{0.920} & \textbf{$-$0.023} & SPI는 회복처럼, SPI(t)는 정체 — 수렴 착시 \\
\end{tbl}
\begin{itemize}
\item 9월 SPI는 8월보다 0.030 상승, SPI(t)는 0.006 상승 — 지수로는 회복, 시간으로는 정체
\item 4월의 역방향(+0.103)은 PV 봉우리 — 돈으로 큰 조달 이연이 시간으로는 0.45개월
\end{itemize}
\keymsg{진척 50\%에서 이미 0.023 벌어졌다 — 후반부에는 더 벌어진다. 일정 보고는 SPI(t)·ES·CPM TF로}
\note{§2.4 월별 표(3월·6월 행은 화면에서 생략, 교재 표에 있다). 두 방향의 차이를 다 보인다 — 봄에는 SPI(t)가 관대(ES가 PV 완만 구간에 있어서), 여름부터는 SPI가 관대(수렴). 즉 어느 쪽이 항상 보수적인 것이 아니라 PV 곡선의 모양과 진척 단계에 따라 다르고, 그래서 둘을 병기한다. 5월 SPI 회복(0.784 \ra{} 0.871)이 라이선스 인도(EV +6.0)로 생긴 것이지 개발 회복이 아니라는 점은 열별 분해(1교시)와 연결된다. 예상 소요 4분.}
\end{frame}

\begin{frame}{Earned Schedule 그림 — EV 73.43이 PV 곡선에서 도달한 시점}
\fig[0.58]{fig_2_4_earned_schedule}
\figcap{그림 2-4. 위 — EV 73.43이 PV 곡선에서 도달한 시점(ES 8.28개월)과 경과 기간(AT 9)의 간격이 SV(t) $-$0.72개월. 아래 — SPI는 종료로 갈수록 1로 수렴하고 SPI(t)는 그렇지 않다. 온담에서 둘은 7월부터 벌어지기 시작했다. Lipke(2003)를 워크북 ⑫ 항목 6의 값으로 재구성}
\keymsg{위 그림의 수평 간격이 일정 편차다 — 수직 간격(SV $-$4.40)은 돈이지 시간이 아니다}
\note{§2.4 그림 2-4. 위 패널에서 EV 수평선이 PV 곡선과 만나는 점을 x축으로 내려 8.28을 읽는 동작을 레이저로 보인다 — SV(수직)와 SV(t)(수평)의 차이가 한 그림에 있다. 아래 패널의 두 선이 7월부터 갈라지는 자리를 짚고, 앞으로 갈수록 위쪽 선(SPI)만 1로 올라갈 것임을 손으로 그려 준다. 예상 소요 3분.}
\end{frame}

\begin{frame}{CPI 안정성 — Christensen 규칙이 말하는 것과 말하지 않는 것}
\begin{itemize}
\item 연구: 미 국방 획득 사업 대규모 표본에서 누적 CPI가 \textbf{진척 20\% 이후 ±10\% 안에서 안정} — "CPI는 대개 좋아지지 않는다"
\item 말하지 않는 것 ① \textbf{표본이 다르다} — 원가 정산 계약 vs 온담 고정가 SI(그 열의 CPI는 계약이 만든 1)
\item ② \textbf{안정된 값이 무엇인지가 먼저} — 온담 누적 CPI 0.978 \ra{} 0.994 \ra{} \textbf{1.073(5월)} \ra{} 1.032은 03-13 지급 조건 변경의 흔적. 조정 0.968이 안정값
\item ③ \textbf{진척 50\%는 20\%가 아니다} — 규칙이 성립한다면 0.968은 크게 회복되지 않는다 \ra{} PMB 초과(TCPI 1.034)는 확정
\item 착시 포함 CPI에 규칙을 적용하면 EAC1 141.23 — PMB보다 낮은, 고정가 구조에서 불가능한 숫자
\end{itemize}
\keymsg{Christensen 규칙은 "CPI가 좋아질 것"이라는 희망을 꺾는 데 쓰는 것이지, 착시 포함 CPI를 외삽하는 데 쓰는 것이 아니다}
\note{§2.5. Christensen \& Payne(1992, Journal of Parametrics)·Christensen \& Heise(1993, NCMJ) — 서지와 후속 논쟁은 6.3절(마무리 교시). 규칙의 실무 함의는 명확하다("20\%에서 CPI 0.9면 곧 회복된다는 말을 믿지 말라")지만 표본·계약 구조·착시 조정 세 조건을 확인하고 쓴다. 온담 5월의 1.073 점프를 학생에게 "원가 효율이 좋아진 것인가"로 물으면 대부분 그렇다고 답한다 — 착시 1의 예고편. 예상 소요 4분.}
\end{frame}

\begin{frame}{흐름 지표 — EVM이 못 보는 것: throughput · WIP · work item age}
\fig[0.50]{../그림/PM_Day3/fig_2_7_flow_metrics.pdf}
\figcap{그림 2-7. S7\til{}S18 throughput(FP/스프린트) 620·620·600·480·470·560·560·600·620·620·640·630, WIP 최고 41(S10), work item age 최장 12영업일}
\keymsg{S10\til{}S11의 470\til{}480은 SPI에 두 달 뒤 나타났다 — 흐름 지표는 EVM의 선행 지표다}
\note{교재 2.7 "흐름 지표", 워크북 §2.5 항목 9. Vacanti의 네 지표를 온담 스프린트 데이터로 읽는다. throughput 하락(S10 480·S11 470)은 R-05 설계 지연이 구현 스프린트로 번진 시점이고, 같은 사실이 월별 SPI에는 5\til{}6월(0.871·0.881)에 완만하게만 나타난다 — 스프린트 단위 지표가 월 단위 EVM보다 먼저 움직인다. WIP 41(S10)은 설계 승인 대기가 쌓인 것이고, work item age 12일은 "완료된 것만 재는" cycle time에서는 보이지 않는다. "완료 기준 cycle time만 개선되는 착시"와 SPI 수렴 착시는 같은 병리라는 점을 학생에게 묻는다(둘 다 끝난 것만 세기 때문이다). CFD는 그림의 하단 — 진행 중 띠의 두께가 WIP다. 예상 소요 5분.}
\end{frame}

\begin{frame}{DORA 4지표의 위치 — 배포 파이프라인의 건강, 프로젝트 성과와는 다른 층}
\begin{itemize}
\item 온담 P1 [추가 설정] — 배포 빈도 주 2회 · 변경 리드타임 3.1일 · 변경 실패율 12\% · 복구 4시간
\item DORA는 \textbf{릴리스 파이프라인}을 재고, EVM은 \textbf{기준선 대비 편차}를 잰다 — 둘은 대체 관계가 아니다
\item 실패율 12\%는 8월 시험 정지(TC 210 \ra{} 95)와 같은 시기 — 시험이 줄면 배포 실패가 늘어난다
\item 하이브리드의 세 층: 스프린트(흐름·DORA) / 월(EVM·ES) / 게이트(허용범위·예외)
\item 한 층의 지표를 다른 층에 올리지 않는다 — 스토리포인트는 EV가 아니고, 배포 빈도는 진척이 아니다
\end{itemize}
\keymsg{세 층은 각자 다른 질문에 답한다 — "잘 흐르는가 / 계획대로인가 / 허용범위 안인가"}
\note{교재 2.7 후반. DORA 값은 회사 정본에 없어 [추가 설정]으로 워크북 §2.5 항목 9에 두었다. 강조할 것은 층의 분리다 — 현장에서 가장 흔한 혼동은 velocity(스토리포인트/스프린트)를 진척률로 보고하는 것이다. 온담은 FP 인수 기준(Rules of Credit)으로 EV를 재므로 velocity와 EV가 분리되어 있고, 그것이 이 프로젝트 성과 측정의 강점이다. 변경 실패율 12\%가 8월 시험 인력 이동과 겹친다는 관찰은 5교시 결함 곡선으로 이어진다. 예상 소요 3분.}
\end{frame}

\begin{frame}{임계경로 — A10 $-$10영업일이 컷오버 창을 밀어낸다}
\fig[0.50]{../그림/PM_Day3/fig_2_8_cpm_asis_tobe.pdf}
\figcap{그림 2-8. as-is: A10 착수 09-28(기준선 180 \ra{} 실제 190) \ra{} MG3 2027-03-05 \ra{} 1차 컷오버 창(02-26) 이탈 \ra{} 2차 창 03-12\til{}14 \ra{} G4 06-11(+10). 합류 버퍼 15 전부 소진}
\keymsg{일정 편차는 돈이 아니라 달력에 나타난다 — SV $-$4.40은 "컷오버 창을 놓친다"는 문장으로 번역되어야 결정이 된다}
\note{교재 2.8 "임계경로 소진", 워크북 §2.5 항목 10. Day2 CPM 표의 A10(구현 R3, 기준선 착수 작업일 180)이 R2 종료 지연(+10)으로 190에 착수했고 임계경로라 그대로 MG3(G3 게이트)를 10영업일 민다. 1차 컷오버 창 02-26\til{}28은 주말 고정이므로 다음 창 03-12\til{}14로 넘어가고, 이것이 R-14(두 번 컷오버)·R-17(2027 자금 집중)을 만든 원인이다. Day2가 남겨 둔 합류 버퍼 15영업일은 여기서 전부 쓰인다 — "버퍼는 누구의 것인가"(Day2 §3 여유 소유권)가 실전이 되는 순간. 학생에게 as-is 그림만 보이고 "무엇을 결정해야 하는가"를 먼저 묻고 다음 프레임의 3대안으로 간다. 예상 소요 5분.}
\end{frame}

\begin{frame}{회복 계획 3대안과 예외 보고 EX-01 — 10-02 제출 · 10-16 결정 시한}
\begin{tbl}[\scriptsize]{L{0.6cm}L{4.6cm}L{2.6cm}L{3.4cm}}
 & 내용 & 효과 & 대가·조건 \\ \midrule
1 & A14 $-$5(SS 겹침) + A16 $-$5(UAT 병행) + R2 이월 420 \ra{} R3 120·A12 300 + CR-08 A12 편입 & MG3 02-19 회복, 여유 0 & A12 TF 10 \ra{} 0 — 제2 임계경로 생성. \textbf{권고} \\
2 & 2차 컷오버 창 03-12\til{}14 수용 & 압축 없음 & 관리예비비 +1.6(두 번 컷오버 운영비) \\
3 & 옴니 "매장 간 이동 요청" 약 260FP \ra{} R5 이연 & A10 부하 감소 & 범위 축소 — 헌장 §5 문언 · 조건부(11-13 S22 속도 620 미만 시 발동) \\
\end{tbl}
\keymsg{예외 보고는 사고 보고가 아니다 — "허용범위 안에서는 보고하지 않을 권리"의 반대급부이며, 대안과 권고를 실어야 문서가 된다}
\note{교재 2.8 후반과 1.3(예외에 의한 관리), 워크북 §2.5 항목 10\til{}11. EX-01의 구조 — 편차(임계 $-$10, IEAC(t) 18.48 vs 허용 종료 06-18) · 원인(R-05 실현·R2 이월 420FP) · 대안 3 · 권고(대안 1 + 대안 3 조건부) · 결정 시한(10-16, 스폰서). 대안 1의 대가가 "A12의 여유 10을 0으로 만든다"는 점을 학생이 스스로 찾게 한다 — 압축은 항상 다른 경로의 여유를 먹는다. 대안 2는 압축이 아니라 수용이고 대가가 관리예비비(스폰서 결재)라 PM 권한 밖이다. 대안 3은 범위 축소이므로 헌장 문언 변경 = CCB 필수(3교시 규칙 ④). 10-16 결정 시한은 S20 착수(10-12) 뒤라는 점 — 결정이 늦으면 대안 1의 A14 SS 겹침이 불가능해진다. 예상 소요 6분.}
\end{frame}

\begin{frame}{2교시 핵심 정리}
\begin{enumerate}
\item EAC 4공식은 네 개의 \textbf{전제}다 — 스폰서가 기억할 숫자는 둘(EAC4 150.86 · EAC3′ 155.04, 경고 대역)
\item Earned Schedule — ES 8.28 · SPI(t) 0.920 · IEAC(t) 18.48개월 \ra{} 2027-07-12 > 허용 06-18. SPI 0.943은 수렴 착시의 시작
\item 착시 셋을 걷어내면 CPI 1.032 \ra{} 0.968 — "원가가 남는다"는 라이선스 이연이 만든 문장
\item 흐름 지표는 EVM의 선행 지표 — S10\til{}S11 throughput 470이 두 달 뒤 SPI에 나타났다
\item 편차는 결정 요청으로 끝나야 한다 — EX-01, 대안 3, 권고 1, 시한 10-16
\end{enumerate}
\keymsg{3교시는 "그 결정을 문서로 어떻게 통과시키는가" — 변경요청서·6축·CCB·기준선 갱신}
\note{교재 제2장 핵심 정리 3\til{}8. 생각해 볼 질문(제2장 2 — "여러분 회사의 마지막 성과보고에서 SPI(t)를 계산하면 어떻게 되는가")을 던지고 답은 받지 않는다. 오후 4교시에서 이 숫자들을 손으로 다시 만든다는 것을 예고. 예상 소요 4분.}
\end{frame}
